\chapter{HTTP Status Codes Reference}

\section{Why Status Codes Matter}

The HTTP status code tells the frontend whether a request succeeded, failed
due to client error, or failed on the server, before it even reads the
response body. Returning \texttt{200 OK} for everything forces every caller
to re-implement its own error detection, whereas consistent codes make
\texttt{axios} interceptors and global error handlers (Chapter~17) possible.

\section{Reference Table}

\begin{table}[h]
\centering
\begin{tabularx}{\textwidth}{l l X}
\toprule
\textbf{Code} & \textbf{Meaning} & \textbf{Typical MERN Usage} \\
\midrule
200 & OK & Successful \texttt{GET}. \\
201 & Created & Resource created via \texttt{POST}. \\
400 & Bad Request & Malformed request or bad JSON. \\
401 & Unauthorized & Missing/invalid/expired auth. \\
403 & Forbidden & Authenticated but not permitted. \\
404 & Not Found & Resource doesn't exist. \\
409 & Conflict & Duplicate data (e.g.\ email taken). \\
422 & Unprocessable Entity & Fails validation rules. \\
500 & Internal Server Error & Unexpected backend failure. \\
\bottomrule
\end{tabularx}
\end{table}

\begin{notebox}[NOTE]
Some teams use \texttt{400} for all validation failures instead of
\texttt{422}; this book uses \texttt{400} for malformed requests and
\texttt{422} for schema/validation errors.
\end{notebox}

\begin{warnbox}[WARNING]
Never return \texttt{200 OK} with an \texttt{\{ "error": "..." \}} body ---
it breaks frontend code relying on \texttt{response.ok} or axios's
non-2xx rejection.
\end{warnbox}
